problem:  
Let \( (M^4, g_0) \) be a smooth closed 4‑manifold with positive isotropic curvature.  
Suppose that for each \( \delta_k \to 0 \) one has a well‑defined **Ricci flow with surgery** \( (M, g_k(t)) \) starting from \( g_0 \), non‑collapsing for all \(t>0\), and existing for all \(t\ge 0\).  

For each fixed \(t>0\), let \( \mu_{k,t} := \mathrm{Vol}_{g_k(t)} \) denote the Riemannian volume measure on \( (M,g_k(t)) \).  
Assume that along a subsequence,  
\[
(M, g_k(t), \mu_{k,t}) \to (X_t, d_t, \mu_t)
\]
in the **measured Gromov–Hausdorff** sense for every \(t>0\), with \( (X_t, d_t) \) a compact 4‑dimensional metric space.  

**Question.** — Under these hypotheses, is it true that for each \(t>0\) the limit triple \( (X_t, d_t, \mu_t) \) satisfies a synthetic lower Ricci curvature bound in the sense of \(\mathrm{RCD}(K(t),4)\)?  
Equivalently, is the long‑time limit of 4‑dimensional Ricci flows with vanishing surgery scale canonically realized, at each time slice, as a (possibly time‑dependent) RCD space?  
If not, can one identify geometric obstructions preventing the limit from satisfying any synthetic curvature bound?

Why it is a "good" problem:  
This refined problem clearly isolates a **precise analytic question** at the intersection of Ricci flow and synthetic curvature geometry: whether the metric‑measure limits \((X_t, d_t, \mu_t)\) arising from 4‑dimensional surgery sequences inherit curvature‑dimension bounds in the sense of \(\mathrm{RCD}(K,N)\) theory.  
It explicitly specifies the convergence mode (measured Gromov–Hausdorff) and the measure (\(\mu_t\)), removing earlier ambiguity.  
The problem’s central simplicity—“Are surgery limits RCD spaces?”—masks exceptional depth: it asks whether the analytic smoothness of Ricci flow transitions naturally into the synthetic metric framework when surgeries vanish.  
A positive resolution would bridge **Perelman’s smooth analytic theory** and **synthetic Ricci‑curvature geometry**, suggesting a canonical, weak formulation of Ricci flow beyond singularities.  
A negative resolution would expose geometric mechanisms (e.g. collapsed necks, measure loss, or topological disconnection) preventing synthetic curvature preservation.  
Either outcome would profoundly influence 4‑manifold geometry and the development of weak Ricci flow theories, making this a conceptually clear yet deeply influential research‑level problem.